\begin{ex}
Hai đường thẳng song song có bao giờ cắt nhau không? (Có ghi 1, không ghi 0).
\shortans{0}
\loigiai{
Hai đường thẳng song song không bao giờ cắt nhau (0).
}
\end{ex}

\begin{ex}
Hai đường thẳng song song có bao nhiêu điểm chung?
\shortans{0}
\loigiai{
Không có điểm chung nào (0).
}
\end{ex}

\begin{ex}
Trong hình chữ nhật ABCD, cạnh AB song song với cạnh nào? (Cạnh CD ghi 1, cạnh BC ghi 2).
\shortans{1}
\loigiai{
Cạnh CD (1).
}
\end{ex}

\begin{ex}
Trong hình chữ nhật ABCD, cạnh AD song song với cạnh nào? (Cạnh BC ghi 1, cạnh CD ghi 2).
\shortans{1}
\loigiai{
Cạnh BC (1).
}
\end{ex}

\begin{ex}
Hình chữ nhật có mấy cặp cạnh đối diện song song với nhau?
\shortans{2}
\loigiai{
Có 2 cặp cạnh đối diện song song.
}
\end{ex}

\begin{ex}
Hình vuông có mấy cặp cạnh đối diện song song với nhau?
\shortans{2}
\loigiai{
Có 2 cặp cạnh đối diện song song.
}
\end{ex}

\begin{ex}
Hình bình hành có mấy cặp cạnh đối diện song song với nhau?
\shortans{2}
\loigiai{
Có 2 cặp cạnh đối diện song song.
}
\end{ex}

\begin{ex}
Hình thoi có mấy cặp cạnh đối diện song song với nhau?
\shortans{2}
\loigiai{
Có 2 cặp cạnh đối diện song song.
}
\end{ex}

\begin{ex}
Hình thang có mấy cặp cạnh đáy song song với nhau?
\shortans{1}
\loigiai{
Hình thang có 1 cặp cạnh đáy song song.
}
\end{ex}

\begin{ex}
Hai đường ray xe lửa là hình ảnh của hai đường thẳng gì? (Song song ghi 1, vuông góc ghi 2).
\shortans{1}
\loigiai{
Hai đường thẳng song song (1).
}
\end{ex}

\begin{ex}
Hai cạnh đối diện của một quyển sách giáo khoa là hình ảnh của hai đường thẳng gì? (Song song ghi 1, cắt nhau ghi 2).
\shortans{1}
\loigiai{
Song song (1).
}
\end{ex}

\begin{ex}
Các dòng kẻ ngang trong cuốn vở ô li là hình ảnh của các đường thẳng gì? (Song song ghi 1, vuông góc ghi 2).
\shortans{1}
\loigiai{
Song song (1).
}
\end{ex}

\begin{ex}
Qua một điểm ở ngoài một đường thẳng, có thể vẽ được bao nhiêu đường thẳng song song với đường thẳng đó?
\shortans{1}
\loigiai{
Chỉ vẽ được 1 đường thẳng song song duy nhất.
}
\end{ex}

\begin{ex}
Khoảng cách giữa hai đường thẳng song song có thay đổi không? (Có thay đổi ghi 1, không đổi ghi 0).
\shortans{0}
\loigiai{
Khoảng cách luôn không đổi (0).
}
\end{ex}

\begin{ex}
Nếu đường thẳng a song song với b, và b song song với c thì a có song song với c không? (Có ghi 1, không ghi 0).
\shortans{1}
\loigiai{
Có song song (1).
}
\end{ex}

\begin{ex}
Nếu đường thẳng a vuông góc với c, và b vuông góc với c thì a và b thế nào với nhau? (Song song ghi 1, vuông góc ghi 2).
\shortans{1}
\loigiai{
Song song với nhau (1).
}
\end{ex}

\begin{ex}
Trong hình vuông ABCD, cạnh AB có song song với cạnh BC không? (Có ghi 1, không ghi 0).
\shortans{0}
\loigiai{
Không, AB vuông góc với BC.
}
\end{ex}

\begin{ex}
Trong tam giác ABC, có cặp cạnh nào song song với nhau không? (Có ghi 1, không ghi 0).
\shortans{0}
\loigiai{
Tam giác không có cạnh nào song song với nhau (0).
}
\end{ex}

\begin{ex}
Hình lục giác đều có bao nhiêu cặp cạnh đối diện song song với nhau?
\shortans{3}
\loigiai{
Có 3 cặp cạnh đối song song.
}
\end{ex}

\begin{ex}
Hai thanh đứng của một chiếc thang là hình ảnh của hai đường thẳng gì? (Song song ghi 1, vuông góc ghi 2).
\shortans{1}
\loigiai{
Song song (1).
}
\end{ex}

\begin{ex}
Các bậc thang ngang của một chiếc thang có song song với nhau không? (Có ghi 1, không ghi 0).
\shortans{1}
\loigiai{
Có song song (1).
}
\end{ex}

\begin{ex}
Bậc thang ngang có vuông góc với thanh đứng của thang không? (Có ghi 1, không ghi 0).
\shortans{1}
\loigiai{
Có vuông góc (1).
}
\end{ex}

\begin{ex}
Có bao nhiêu cặp cạnh song song trong một hình chữ nhật?
\shortans{2}
\loigiai{
Có 2 cặp.
}
\end{ex}

\begin{ex}
Hình bình hành có bao nhiêu cạnh?
\shortans{4}
\loigiai{
Có 4 cạnh.
}
\end{ex}

\begin{ex}
Hình bình hành có bao nhiêu đỉnh?
\shortans{4}
\loigiai{
Có 4 đỉnh.
}
\end{ex}
